\section{Cell weights and three-statistic disc detection}\label{sec:geometry}

\subsection{Validated finite geometry}

The existing validator checks reciprocal face pairings, face-connectedness,
ambient orientability, absence of reversed edge identifications, spherical
or disc vertex links, and a single finite boundary torus. Ideal vertices and
other boundary topologies are rejected. Normal coordinates are checked for
nonnegativity, face matching, and the quadrilateral constraints. These checks
are retained before any new component query.

The proof requires no irreducibility hypothesis on the supplied manifold.
Irreducibility becomes relevant to other recognition and hierarchy operations,
not to the classification of its embedded surface components. In particular,
the implementation must correctly reject a closed projective-plane component
as a disc even though its Euler characteristic is one.

For every global triangulation edge $e$, let $w_e$ be its normal intersection
number. A block of $w_e$ consecutive integers represents those intersection
points, ordered using a consistent orientation of that edge. Concatenate
these blocks to obtain a universe of size
\[
                        N=\sum_e w_e.
\]
Each normal arc in a face pairs its two endpoint positions. Consecutive arcs
of a fixed type give one interval isometry. Count a glued face once and each
boundary face once. The number of pairings is at most $12t$. The graph formed
by these vertices and arc edges is the one-skeleton of the normal surface's
polygonal cellulation. Its connected components are exactly those of the
surface. Identifications within one tetrahedron, loop arcs, and multiple
edges cause no problem for either connectivity or the Euler formula.

\subsection{Integral ownership of Euler characteristic}

One can distribute fractional Euler weights over polygon corners. We instead
use an integral ownership rule, which avoids denominators entirely.
Give every edge-intersection point an initial Euler weight $+1$. For every
normal arc, charge $-1$ to one of its endpoints. For every normal polygon,
charge $+1$ to one selected corner. The selected corner belongs to the same
surface component as the entire polygon; likewise the endpoint of an arc.
Therefore, for every component $C$,
\[
             \sum_{p\in C} w_\chi(p)=V(C)-E(C)+F(C)=\chi(C).
\]
This is true even when several selected local corners become the same global
intersection point: their additive contributions must then accumulate there.

All ownership data is described by $O(t)$ intervals. The negative arc charge
is assigned on the canonical source interval of each maintained pairing.
For a triangle type at vertex $v$, choose a fixed incident local edge and use
the triangle points nearest $v$. For an active quadrilateral type, choose the
lexicographically first local edge that it crosses. Along an edge from local
vertex $a$ to local vertex $b$, the points appear as triangles at $a$, then
quadrilaterals, then triangles at $b$. If their relevant counts are $T_a,Q$,
the selected quadrilateral corners occupy
\[
                            [T_a,T_a+Q).
\]
Reflect this interval if the local edge direction disagrees with the global
orientation. No stack is expanded.

\begin{lemma}[Cell-weight identity]\label{lem:cell-weight}
The constructed piecewise-constant integer weight has component sums
$\chi(C)$, and has an $O(t)$-interval additive description. A second weight
obtained by retaining just the polygon credits sums to the number of normal
polygons in $C$.
\end{lemma}
\begin{proof}
Every cell receives exactly one charge of its appropriate sign and each
charge lies in its own connected component. A tetrahedron has at most five
nonempty polygon types, and a face has at most three arc types. The listed
charges are constant on each type's consecutive set of corners or endpoints.
\end{proof}

The full-coordinate reference replaces each polygon credit by the corresponding
standard basis vector of $\Z^{7t}$. A component's orbit sum is then its actual
normal coordinate vector. This is the classical component-recovery route
\cite{AHT2006}; the new disc query avoids its growing vector dimension.

\subsection{A small, certified homology basis of the boundary}

Let $G$ be the primal one-skeleton of the boundary triangulation. Self-loops
and multiple edges are permitted. Select a spanning tree $T$ of $G$. In the
dual graph select a spanning tree $T^*$ using edges whose primal edges are
outside $T$. Exactly two primal edges are in neither tree, since the boundary
is a torus and
\[
 |E|-(|V|-1)-(|F|-1)=2.
\]
For each leftover edge, append the unique path in $T$ connecting its endpoints.
These two fundamental cycles, $c_1,c_2$, form a basis of
$H_1(\partial\mathcal T;\F_2)$.

For completeness, the tree-cotree fact can be seen by contracting the primal
tree and removing edges dual to a dual spanning tree. These operations
preserve the surface and reduce its cellular description to one vertex,
one face, and two edges. Over $\F_2$, the boundary of that face is zero,
and the two surviving cycles freely generate the two-dimensional first
homology. Reversing the contractions gives exactly the chosen fundamental
cycles. The dual graph outside the primal tree is connected: a separating
dual cut would correspond to a nonempty primal cycle contained in the tree.

The certificate checker does not rely on the spanning-tree algorithm's output
being correct. It independently checks that each proposed cycle has even
incidence at every vertex, constructs the boundary vectors of all triangles,
and performs binary Gaussian elimination. The two proposed cycles must each
increase the rank modulo the face-boundary span. Since the validator has
already established a torus, two independent classes are a basis. Repeated
edges of a triangular face cancel modulo two in this check.

\subsection{Boundary intersection weights}

For each $j\in\{1,2\}$, give every intersection point on a boundary edge in
$c_j$ a homology weight $+1$ in coordinate $j$; give every other point weight
zero. If $z_e(C)$ counts the boundary points of component $C$ on edge $e$, then
its orbit sum is
\[
                    \eta_j(C)=\sum_{e\in c_j} z_e(C).
\]
The code retains this as an exact integer and reduces it modulo two when
interpreting the result. Equivalently, one could carry an $\F_2$ coordinate
throughout the weighted algorithm; the present implementation uses one
uniform integer-vector engine for all query modes.

The parity vector $z_e(C)\bmod2$ is the cellular cocycle Poincare-dual to
$\partial C$: every boundary triangle meets the normal boundary curve in
arcs and therefore has even total endpoint incidence. Evaluating that cocycle
on the two basis cycles yields
\[
    h(C)=(\eta_1(C),\eta_2(C))\bmod2,
    \qquad
    h(C)=0\ \Longleftrightarrow\ [\partial C]=0
       \text{ in }H_1(\partial\mathcal T;\F_2).
\]
The equivalence uses the nondegenerate homology-cohomology pairing over the
field $\F_2$. It applies to the total boundary of each component, rather than
to one arbitrarily selected boundary circle.

\subsection{Three additive weights characterize compressing discs exactly}
\label{sec:three-weight-criterion}

Let $C$ be a connected component of the supplied proper normal surface.
The cell-incidence weight gives $\chi(C)$. Let $\gamma_1,\gamma_2$ be the
checked homology basis of the boundary torus described above, and define
\[
 h_i(C)=\sum_{e\in\gamma_i}|C\cap e|\pmod2,
 \qquad i=1,2.
\]
The raw orbit calculation retains the corresponding integer sums; parity is
taken only when interpreting the resulting component signature. In
particular, signed Euler arithmetic and binary homology arithmetic are not
silently identified.

\begin{lemma}
\label{lem:boundary-character-completeness}
The pair $(h_1(C),h_2(C))$ vanishes exactly when the mod-two boundary class of
$C$ vanishes in the ambient torus.
\end{lemma}
\begin{proof}
On every boundary triangle, a normal arc contributes two endpoints.
Consequently the parity of edge intersections is a cellular $1$-cocycle,
including in generalized triangulations where repeated face edges must be
combined modulo two. This cocycle represents the Poincar\'e dual of the
boundary multicurve. Evaluation on a homology basis determines its cohomology
class, since evaluation gives the nondegenerate pairing between
$H^1(T^2;\F_2)$ and $H_1(T^2;\F_2)$. Both spaces have dimension two. Thus
the two values vanish precisely for the zero boundary class.
\end{proof}

\begin{theorem}[Exact three-weight disc criterion]
\label{thm:three-weight-disc}
A connected component $C$ is a compressing disc for the ambient boundary
torus if and only if
\[
 \chi(C)=1,
 \qquad (h_1(C),h_2(C))\ne(0,0).
\]
\end{theorem}
\begin{proof}
A compressing disc has Euler characteristic one. Its boundary is an embedded
essential circle on the torus. Such a circle represents a primitive integral
slope, so its reduction modulo two is nonzero. The lemma proves necessity.

Conversely, the nonzero boundary class implies that $C$ has boundary. The
classification of connected compact surfaces now leaves only a disc. Indeed,
an orientable surface of genus $a$ with $b\ge1$ boundary circles has
$\chi=2-2a-b$; equality to one forces $a=0,b=1$. A nonorientable surface with
$k\ge1$ crosscaps and $b\ge1$ boundary circles has
$\chi=2-k-b\le0$. The resulting disc has nonzero boundary class, and hence
its simple boundary circle is essential on the torus.
\end{proof}

The boundary condition in this argument is necessary. A closed projective
plane also has Euler characteristic one; its two boundary evaluations are
zero, so it is correctly excluded. Such closed one-sided surfaces occur in
the supported general torus-boundary manifold fixtures, even though those
fixtures need not be knot exteriors. No assumption of irreducibility is
smuggled into the local component test.

Two binary linear probes are also the minimum needed to recognize the zero
element of the full torus boundary-class space by linear evaluation alone.
A single linear functional on the two-dimensional space has a nonzero
kernel. The two-probe statement concerns this unrestricted homology interface;
it does not prove a lower bound for a specially restricted collection of
normal surfaces in one fixed manifold, nor for nonlinear preprocessing.
Together with the Euler coordinate, it explains the natural three-weight
implementation and its reduction from $7t$-dimensional coordinate extraction.

The torus hypothesis cannot be dropped from the equivalence. On a boundary
surface of genus at least two, an essential separating simple circle can have
zero mod-two homology. The same predicate would then give only a sufficient
test for certain compressing discs, and could miss others. The native input
validator's one-torus-boundary contract is therefore part of the theorem,
not merely a convenience for constructing a small basis.

Finally, the test must be applied to individual component weights. Two
parallel meridians have zero total boundary class but contribute two
compressing discs. Conversely, an inessential vertex-link disc together with
a M\"obius band can have total Euler characteristic one and nonzero total
boundary class while contributing no compressing disc. A weighted orbit
histogram resolves both ambiguities without listing the components one by one.

\subsection{A proved two-weight refinement using ambient homology}
\label{sec:two-weight-refinement}

The three-weight implementation uses only the boundary torus when choosing
its homology probes. The ambient manifold supplies additional information:
boundaries of proper surfaces occupy a one-dimensional subspace of the
torus homology. Exploiting that restriction gives the following proved
refinement. It is \emph{not implemented or benchmarked in this delivery};
the reported disc-query improvements use the implemented three-weight mode
and the stated comparison arms.

\begin{samepage}
\begin{proposition}[One ambient-certified boundary probe]
\label{prop:two-weight-disc}
Let $M$ be a compact connected orientable three-manifold with exactly one
torus boundary, and let $i:\partial M\hookrightarrow M$. There exists a
boundary cycle $\gamma$ with $i_*[\gamma]\ne0$ in $H_1(M;\F_2)$. For any
such cycle and every connected proper normal-surface component $C$,
\[
 C\text{ is a compressing disc}
 \quad\Longleftrightarrow\quad
 \chi(C)=1\ \text{ and }\ |\partial C\cap\gamma|\equiv1\pmod2.
\]
The intersection count can be computed by the same additive edge-point
weights as before. Thus two integer weights suffice after the indicated
ambient preprocessing.
\end{proposition}
\end{samepage}
\begin{proof}
All homology groups in this proof have coefficients in $\F_2$. Write
$b_j=\dim H_j(M)$. Poincar\'e--Lefschetz duality gives
$\chi(M,\partial M)=-\chi(M)$; combined with the Euler characteristic of
the pair, this gives $2\chi(M)=\chi(\partial M)=0$. Since $M$ is connected
with nonempty boundary, $b_0=1,b_3=0$, and hence $b_1-b_2=1$
\cite{Hatcher2002}.

In the long exact sequence of the pair, $H_0(\partial M)\to H_0(M)$ is an
isomorphism. Consequently $H_1(M)\to H_1(M,\partial M)$ is surjective.
Duality identifies the dimension of the latter group with $b_2$. Exactness
therefore gives $\dim\operatorname{im}i_*=b_1-b_2=1$. The kernel
\[
 L=\ker\bigl(H_1(\partial M)\longrightarrow H_1(M)\bigr)
\]
is a line in the two-dimensional torus homology. Every compact surface,
including a nonorientable one, has a relative fundamental class over
$\F_2$. Pushing that class into $(M,\partial M)$ shows that
$[\partial C]\in L$.

The torus intersection pairing is nondegenerate and alternating. A line
$L$ consequently has orthogonal complement $L$ itself. Since
$i_*[\gamma]\ne0$, the class $[\gamma]$ is outside $L$ and pairs to one
with its unique nonzero element. Odd intersection with $\gamma$ is therefore
equivalent to $[\partial C]\ne0$. The surface-classification and primitive
torus-slope argument already proved then gives the disc criterion. This
argument imposes neither irreducibility nor a torsion-free integral homology
assumption.
\end{proof}

There is a concrete finite preprocessing algorithm. Embed the two verified
boundary basis cycles into the ambient global edge-chain space over $\F_2$.
Form the span of boundaries of all ambient triangular face classes, counting
paired faces once and combining repeated edges modulo two. At least one
boundary basis cycle lies outside this span, because the boundary inclusion
has rank one. Choose such a cycle as $\gamma$. It is already a cycle, so
being outside the face-boundary span is exactly the required nonzero ambient
homology condition.

The preprocessing can carry a particularly small certificate: an ambient
cellular $1$-cocycle $u$ satisfying
\[
 u(\partial f)=0\quad\text{for every triangular face }f,
 \qquad u(\gamma)=1.
\]
These identities prove that $\gamma$ is not a boundary. A producer can find
$u$ by Gaussian elimination in $O(t^3)$ elementary operations over $\F_2$;
replay checks $O(t)$ constant-size face equations and the cycle evaluation,
without repeating elimination. The cycle and cocycle depend only on the
triangulation and can be reused across its candidate surfaces.

Both signatures have constant dimension. This refinement saves one
transported integer coordinate while adding ambient preprocessing; its practical benefit remains to be measured. It
does not contradict the earlier observation about two linear probes on the
\emph{unrestricted} torus class space. The new probe uses a certified
geometric restriction to the line $L$. With several boundary components the
rank argument above changes, so this particular one-probe conclusion must
not be transferred to that setting without a new analysis.


\subsection{Component histograms and witness extraction}

In disc mode the weighted output is initially a histogram of exact integer
triples. The result interpreter reduces its last two entries modulo two and
merges identical triples, adding their binary multiplicities. This last
aggregation is safe because the predicate depends only on those three
displayed values. It is unnecessary to distinguish normally nonisotopic
components that have the same disc-test signature.

Summary mode adds the number of polygons and the number of boundary
intersection points. The latter equals the number of boundary normal arcs
componentwise, since every boundary circle has an equally sized vertex and
edge set in its induced cellulation. These two counts are diagnostic data;
they do not determine the number of boundary circles or the full genus.

When an actual disc vector is required, the full-coordinate mode recovers a
histogram of connected-component coordinate vectors. One may select any
record with the certified disc predicate. The coordinate vector is explicit
in $7t$ binary entries, while a large number of identical copies is retained
as a multiplicity. The existing unweighted trace can be checked and reused,
so extraction does not require another search for an orbit-reduction order.
This separates a cheap decision query from a more informative extraction
query without conflating their costs.

\subsection{Why the geometric restriction matters}

The boundary-torus assumption is substantive. On a boundary surface of genus
at least two, an essential separating simple closed curve has zero mod-two
homology. Such a curve could bound a compressing disc and would escape the
three-weight test. A higher-genus extension therefore needs compressed
curve-essentiality information, not just more copies of a homology basis.
The supplied validator rejects these inputs, so the implementation does not
silently apply the torus theorem outside its hypotheses.
